\section{Analysis - Folgen und Reihen}
\subsection{Folgen}
\subsubsection{Folgen und Grenzwerte}
\begin{frame}{Folgen}
	\begin{definition}
		Eine \textbf{reellwertige Folge} $(a_{n})_{n \in \mathbb{N}}$ ist eine Abbildung von $\mathbb{N}$ in $\mathbb{R}$, d. h. jedem $n \in \mathbb{N}$ wird eine reelle Zahl $a_{n} \in \mathbb{R}$ zugeordnet.
		
		Falls $a_{n} = c$ für alle $n \in \mathbb{N}$ und eine feste Zahl $c$, so heißt die Folge $(a_{n})_{n \in \mathbb{N}}$ \textbf{konstante Folge}.
		
		Eine Folge, für die $a_{n+1} = q \cdot a_{n}$ erfüllt ist für alle $n \in \mathbb{N}$, heißt \textbf{geometrische Folge}.
	\end{definition}
	\begin{example}
		Beispiele für Folgen sind:
		\begin{align*}
			(a_{n})_{n \in \mathbb{N}} = (n)_{n \in \mathbb{N}} & \mbox{ also } (1, 2, 3, 4, 5, 6, \ldots) \\
			(a_{n})_{n \in \mathbb{N}} = \left(\frac{1}{n}\right)_{n \in \mathbb{N}} & \mbox{ also } (1, \frac{1}{2}, \frac{1}{3}, \frac{1}{4}, \frac{1}{5}, \frac{1}{6}, \ldots) \\
			(a_{n})_{n \in \mathbb{N}} = \left((-1)^{n}\frac{1}{n}\right)_{n \in \mathbb{N}} & \mbox{ also } (-1, \frac{1}{2}, -\frac{1}{3}, \frac{1}{4}, -\frac{1}{5}, \frac{1}{6}, \ldots)	
		\end{align*}
	\end{example}
\end{frame}

\begin{frame}{Beispiele für Folgen}
	\begin{definition}
		Die Folge
		\[
		a_{n} := \frac{1}{n} \mbox{ mit } \left(a_{n}\right) = \left(\frac{1}{1}, \frac{1}{2}, \frac{1}{3}, \frac{1}{4}, \ldots \right)
		\] heißt \textbf{harmonische Folge}.
	\end{definition}
	
\end{frame}

\begin{frame}{Beispiele rekursiv definierter Folgen}
		Man kann folgen auch \textbf{rekursiv} definieren. 
	\begin{example}
		Die \textbf{Fibonacci-Folge} wird definiert als:
		\[
			a_{1} := 1 \qquad a_{2} := 1 \qquad a_{n} := a_{n-1} + a_{n-2} \mbox{ für } n \geq 3
		\]
		also 
		\[
			 \left( a_{n} \right) = \left( 1, 1, 2, 3, 5, 8, 13, 21, \ldots \right) 
		\]
	\end{example}
	\begin{example}
		Für die Folge
		\[
		a_{1} := 1 \qquad a_{n} := n \cdot a_{n-1} \mbox{ für } n \geq 2
		\]
		ist die Folge der Fakultäten:
		\[
			a_{n} = n!
		\]
	\end{example}
\end{frame}

\begin{frame}{Numerische Berechnung von Wurzeln}
	\begin{example}
		Für eine Zahl $x > 0$ sei $a_{1} := 1$ und
		\[
			a_{n+1} := \frac{1}{2} \left( a_{n} + \frac{x}{a_{n}} \right) \mbox{ für alle } n \geq 1 
		\]
		Diese Folge nennen wir hier \textbf{Heron-Folge zur Zahl $x$}.
	\end{example} 
\end{frame}

\begin{frame}{Folgen}
	\begin{bemerkung}
		Es gibt auch allgemeinere Folgen, die jedem $n \in \mathbb{N}$ ein Element einer beliebigen Menge $M$ zuordnet.
		
		Die Zahlen bzw. Elemente $a_{n}$ heißen \textbf{Glieder der Folge}.
		
		Folgen können beliebige Abfolgen von Elementen der Bildmenge $M$ (hier immer $\mathbb{R}$) sein. Falls es aber eine Regel gibt, nach der die Folgenglieder berechnet werden, nennt man diese Regel das \textbf{Bildungsgesetz der Folge}.
	\end{bemerkung}
\end{frame}

\begin{frame}{Eigenschaften von Zahlenfolgen}
	\begin{definition}
		Eine Folge $(a_{n})_{n \in \mathbb{N}}$ heißt
		\begin{itemize}
			\item \textbf{nach oben beschränkt}, wenn es eine Zahl $M$ gibt, sodass $a_{n} < M$ für alle $n \in \mathbb{N}$.
			\item \textbf{nach unten beschränkt}, wenn es eine Zahl $m$ gibt, sodass $m < a_{n}$ für alle $n \in \mathbb{N}$.
			\item \textbf{beschränkt}, wenn sie nach oben und nach unten beschränkt ist.
			\item \textbf{unbeschränkt}, wenn sie nicht beschränkt ist.
			\item \textbf{streng monoton steigend}, wenn $a_{n+1}>a_{n}$ für alle $n \in \mathbb{N}$.
			\item \textbf{monoton steigend}, wenn $a_{n+1}\geq a_{n}$ für alle $n \in \mathbb{N}$.
			\item \textbf{streng monoton fallend}, wenn $a_{n+1}<a_{n}$ für alle $n \in \mathbb{N}$.
			\item \textbf{monoton fallend}, wenn $a_{n+1}\leq a_{n}$ für alle $n \in \mathbb{N}$.
			\item \textbf{alternierend}, wenn $a_{n} \cdot a_{n+1} < 0$ für alle $n \in \mathbb{N}$, d. h. die Vorzeichen benachbarter Folgenglieder wechseln.
		\end{itemize}
	\end{definition}
\end{frame}

\begin{frame}{Eigenschaften von Zahlenfolgen}
	\begin{definition}
		Eine Zahl $h$ heißt \textbf{Häufungspunkt der Folge} $(a_{n})_{n \in \mathbb{N}}$, wenn für jedes $\varepsilon>0$ 
		\[
			\left| a_{n} - h \right| < \varepsilon
		\]
		für \textbf{unendlich} viele Folgenglieder gilt.
		
		Das heißt anschaulich: Egal wie nah man $h$ kommt, es gibt immer noch unendlich viele Folgenglieder, die noch näher an $h$ liegen.
		
		Hat eine Folge \textbf{genau einen} Häufungspunkt $h$, so nennt man die Folge \textbf{konvergent} und $h$ den \textbf{Grenzwert} der Folge und man schreibt:
		\[
			\lim_{n \to \infty} a_{n} = h
		\]
		
		Eine \textbf{nicht konvergente} Folge heißt \textbf{divergent}.
	\end{definition}
\end{frame}

\begin{frame}{Konvergenz von Folgen}
	\begin{theorem}
		Eine Folge ist genau dann kovergent gegen $h \in \mathbb{R}$, wenn es zu jedem $\varepsilon > 0$ eine Zahl $N(\varepsilon) \in \mathbb{N}$ gibt, sodass
		\[
			\left| a_{n} - h \right| < \varepsilon \qquad \forall n > N(\varepsilon)
		\]  
	\end{theorem}
	\begin{bemerkung}
		Eine konvergente Folge, deren Grenzwert $h = 0$ ist, heißt \textbf{Nullfolge}.
		
		Wächst eine Zahlenfolge über alle Grenzen, so schreibt man:
		\[
		\lim_{n \to \infty} a_{n} = \infty
		\]
	\end{bemerkung}
\end{frame}

\begin{frame}{Konvergenz und Monotonie}
	\begin{theorem}
		Jede \textbf{beschränkte und monotone} Zahlenfolge ist \textbf{konvergent}.
	\end{theorem}
	\begin{bemerkung}
		Beide Bedingungen (Beschränktheit und Monotonie) sind wichtig, eine allein reicht (noch) nicht aus für die Konvergenz einer Folge.
	\end{bemerkung}
\end{frame}

\begin{frame}{Rechenregeln für Grenzwerte}
	\begin{theorem}
		Es gilt für zwei \textbf{konvergente} Folgen $(a_{n})_{n \in \mathbb{N}}$ und $(b_{n})_{n \in \mathbb{N}}$ mit 
		\[
		\lim_{n \to \infty} a_{n} = a \mbox{  und  } \lim_{n \to \infty} b_{n} = b
		\] 
		\begin{itemize}
			\item Die \textbf{Summenfolge} $(c_{n})_{n \in \mathbb{N}}$ mit $c_{n} := a_{n} + b_{n}$ konvergiert und:
			\[
			\lim_{n \to \infty} c_{n} = \lim_{n \to \infty} a_{n} + b_{n} = \lim_{n \to \infty} a_{n} + \lim_{n \to \infty} b_{n} = a + b
			\]
			\item Die Folge $d_{n} := c \cdot a_{n}$ mit einer Konstanten $c \in \mathbb{R}$  konvergiert:
			\[
			\lim_{n \to \infty} d_{n} = \lim_{n \to \infty} c \cdot a_{n} = c \cdot a
			\]
			\item Die \textbf{Produktfolge} $(p_{n})_{n \in \mathbb{N}}$ mit $p_{n} := a_{n} \cdot b_{n}$ konvergiert mit:
			\[
			\lim_{n \to \infty} p_{n} = \lim_{n \to \infty} a_{n} \cdot b_{n} = \lim_{n \to \infty} a_{n} \cdot \lim_{n \to \infty} b_{n} = a \cdot b
			\]
		\end{itemize}
	\end{theorem}
\end{frame}

\begin{frame}{Rechenregeln für Grenzwerte}
	\begin{bemerkung}
		Aus dem vorherigen Satz folgt für zwei \textbf{konvergente} Folgen $(a_{n})_{n \in \mathbb{N}}$ und $(b_{n})_{n \in \mathbb{N}}$ auch:
		\begin{itemize}
			\item Die Differenzfolge konvergiert:
			\[
			\lim_{n \to \infty} (a_{n} - b_{n}) = \lim_{n \to \infty} a_{n} - \lim_{n \to \infty} b_{n} = a - b
			\]
			\item Sind alle $b_{n} \neq 0$ und $b \neq 0$, so konvergiert auch die \textbf{Quotientenfolge}:
			\[
			\lim_{n \to \infty} \frac{a_{n}}{b_{n}} = \frac{\lim_{n \to \infty} a_{n}}{\lim_{n \to \infty} b_{n}} = \frac{a}{b}
			\]
%						
		\end{itemize}
	\end{bemerkung}
\end{frame}

\begin{frame}{Beispiele für die Berechnung von Grenzwerten}
	\begin{example}
	\begin{align*}
		& \lim_{n \to \infty} \frac{2n}{n+1} = \lim_{n \to \infty} \frac{2n+2 - 2}{n+1}  = \lim_{n \to \infty} \frac{2(n+1) - 2}{n+1} \\ & = \lim_{n \to \infty} \left( 2 - \frac{2}{n+1} \right) = \lim_{n \to \infty} 2 - \lim_{n \to \infty} \frac{2}{n+1}  = 2 - 0   = 2 \\
		& \lim_{n \to \infty} \cos \left( 2- \frac{1}{n} \right)  = \cos 2 \\
		& \lim_{n \to \infty} (-1)^{n} \left( 1 + \frac{1}{n}  \right) \mbox{ existiert nicht! Warum?}
	\end{align*}
	\end{example}
\end{frame}

\begin{frame}{Numerische Berechnung von Wurzeln}
	\begin{theorem}
		Die \textbf{Heron-Folge zur Zahl $x>0$} mit $a_{1} = 1$ und
		\[
		a_{n+1} := \frac{1}{2} \left( a_{n} + \frac{x}{a_{n}} \right) \mbox{ für alle } n \geq 1 		
		\]
		konvergiert und für den Grenzwert gilt:
		\[
			a := \lim_{n \to \infty} a_{n+1} = \lim_{n \to \infty} \frac{1}{2} \left( a_{n} + \frac{x}{a_{n}} \right)
		\]
		und daher für $a$ 
		\[
			a = \frac{1}{2} \left( a + \frac{x}{a} \right)
		\]
		Diese Gleichung hat nur die positive Lösung $a=\sqrt{x}$ und daher:
		\[
			\lim_{n \to \infty} a_{n} = \sqrt{x}
		\]
	\end{theorem}
\end{frame}

\subsubsection{Spezielle Folgen}

\begin{frame}{Arithmetische Folgen}
	\begin{definition}
		Gilt für eine Folge $\left(a_{n}\right)_{n \in \mathbb{N}}$ für alle $n \in \mathbb{N}$:
		\[
		d := a_{n+1} - a_{n} = \mbox{konstant,}
		\]
		so heißt sie \textbf{arithmetische Folge}.
	\end{definition}
\end{frame}	
\begin{frame}{Eigenschaften arithmetischer Folgen}
	\begin{bemerkung}
		\begin{itemize}
		\item Eine arithmetische Folge ist \textbf{eindeutig} durch das Anfangsglied $a_{1}$ und die konstante Differenz $d$ bestimmt und es gilt:
		\[
			a_{n} = a_{1} + d \cdot (n - 1)
		\]
		\item Da für eine arithmetische Folge sowohl
		\[
			d = a_{n+1} - a_{n} \mbox{ als auch } d = a_{n} - a_{n-1}
		\] 
		gilt, ist (durch Gleichsetzen der beiden Gleichungen)
		\[
			a_{n} = \frac{1}{2} \left( a_{n+1} + a_{n-1}\right)
		\]
		also jedes Folgenglied das \textbf{arithemtische Mittel} seiner benachbarten Folgenglieder.
		\end{itemize}	
	\end{bemerkung}
\end{frame}

\begin{frame}{Divergenz/Konvergenz arithmetischer Folgen}
	\begin{theorem}
		Eine \textbf{arithmetische Folge} konvergiert dann und nur dann, wenn $d = a_{n} - a_{n-1} = 0$ ist. In diesem Fall ist die Folge sogar konstant und es gilt:
		\[
		\lim_{n \to \infty} a_{n} = a_{1} = a_{2} = \cdots
		\] 
		
		Gilt für eine arithmetische Folge also $d \neq 0$, dann divergiert die Folge:
		\begin{align*}
			\lim_{n \to \infty} a_{n} & = \infty & \mbox{ falls d > 0} \\
			\lim_{n \to \infty} a_{n} & = - \infty & \mbox{ falls d < 0}
		\end{align*}
	\end{theorem}
\end{frame}

\begin{frame}{Geometrische Folgen}
	\begin{definition}
		Eine Folge $(a_{n})_{n \in \mathbb{N}}$ heißt \textbf{geometrische Folge}, wenn 
		\[
		\frac{a_{n+1}}{a_{n}} = q = \mbox{const. für alle } n \in \mathbb{N}   
		\]
		ist.
	\end{definition}
\end{frame}	
\begin{frame}{Eigenschaften geometrischer Folgen}
	\begin{bemerkung}
		\begin{itemize}
		\item Eine geometrische Folge ist eindeutig durch ihr Anfangsglied $a_{1}$ und den konstanten Quotienten $q$ bestimmt und es gilt:
		\[
			a_{n} = a_{1} \cdot q^{n-1}
		\]
		\item Da für eine geometrische Folge sowohl
		\[
		q = \frac{a_{n+1}}{a_{n}} \mbox{ als auch } q = \frac{a_{n}}{a_{n-1}}
		\]
		gilt, ist  (durch Gleichsetzen der beiden Gleichungen)
		\[
			a_{n} = \sqrt{a_{n-1} \cdot a_{n+1}}
		\]
		also jedes Folgenglied das \textbf{geometrische Mittel} seiner benachbarten Folgenglieder.
		\end{itemize} 
	\end{bemerkung}
\end{frame}

\begin{frame}{Konvergenz/Divergenz geometrischer Folgen}
	\begin{bemerkung}
	Bezüglich der Konvergenz bzw. Divergenz einer geometrischen Folge $a_{n} = a_{1} q^{n-1}$ gilt folgende Fallunterscheidung:
	\begin{itemize}
		\item Falls $|q| < 1$ ist, so konvergiert die Folge gegen 0. 
		\[
		\lim_{n \to \infty} a_{n} = \lim_{n \to \infty} a_{1} q^{n-1} = 0
		\]
		\item Falls $q\leq-1$ ist, so ist die Folge alternierend und divergent.
		\item Falls $q>1$ ist, so ist die Folge divergent.
		\item Falls $q=1$ ist, so konvergiert die Folge, weil sie konstant ist: $\lim_{n \to \infty} a_{n} = \lim_{n \to \infty} a_{1} \cdot 1^{n-1} = a_{1}$.
	\end{itemize}
	\end{bemerkung}	
\end{frame}

\begin{frame}{Techniken zum Nachweis von Beschränktheit und Monotonie}
	\begin{itemize}
		\item Sind alle Folgenglieder positiv bzw. negativ? Falls ja: Die Folge ist durch 0 nach oben bzw. nach unten beschränkt.
		\item Gilt $a_{n+1} - a_{n} \geq 0$ oder $a_{n+1} - a_{n} > 0$? Falls ja: Die Folge ist monoton wachsend bzw. streng monoton wachsend.
		\item Gilt $a_{n+1} - a_{n} \leq 0$ oder $a_{n+1} - a_{n} < 0$? Falls ja: Die Folge ist monoton fallend bzw. streng monoton fallend.
		\item Gilt $\frac{a_{n+1}}{a_{n}} \geq 1$ oder $\frac{a_{n+1}}{a_{n}} > 1$? Falls ja: Die Folge ist monoton wachsend bzw. streng monoton wachsend.
		\item Gilt $\frac{a_{n+1}}{a_{n}} \leq 1$ oder $\frac{a_{n+1}}{a_{n}} < 1$? Falls ja: Die Folge ist monoton fallend bzw. streng monoton fallend.
	\end{itemize}
\end{frame}

\begin{frame}{Berechnung von Grenzwerten - Beispiele}
	\begin{example}
		\begin{align*}
		\lim_{n \to \infty} \frac{2}{n} + \left(\frac{1}{3}\right)^{n} +7 & = 7 \\
		\lim_{n \to \infty} \frac{3n^{2} +7n + 8}{5n^{2} -8n +1} & = \frac{3}{5} \\
		\lim_{n \to \infty} \frac{a_{r} n^{r} + \ldots + a_{1} n + a_{0}}{b_{s} n^{s} + \ldots + b_{1} n + b_{0}} & = \begin{cases}
		0 & \text{falls } r <s \\
		+\infty & \text{falls } s<r \text{ und } \frac{a_{r}}{b_{s}} > 0 \\
		-\infty & \text{falls } s<r \text{ und } \frac{a_{r}}{b_{s}} < 0 \\
		\frac{a_{r}}{b_{s}} & \text{falls s=r} \\
		\end{cases} \\
		\lim_{n \to \infty} \sqrt{n^{2}+3n+3} -n & = \frac{3}{2} \\
		\lim_{n \to \infty} \frac{n}{2^{n}} = 0
 		\end{align*}
	\end{example}
\end{frame} 

\begin{frame}{Grenzwertbestimmung bei rekursiv definierten Folgen}
	\begin{bemerkung}
	Falls eine Folge $(a_{n})$ rekursiv definiert wird (wie z.B. die Heron-Folge), bestimmt man den Grenzwert in der Regel wie folgt:
	\begin{itemize}
		\item Man zeigt, dass die Folge konvergiert z.B. durch
		\begin{itemize}
			\item $(a_{n})$ ist beschränkt und
			\item $(a_{n})$ ist monoton.
		\end{itemize}
		\item Man ersetzt in der Rekursionsvorschrift alle $a_{n+1}$, $a_{n}$ usw. durch den zu berechnenden Grenzwert $a$ und erhält damit eine sogenannte Fixpunktgleichung für a.
		\item Man löst die Fixpunktgleichung aus dem vorherigen Schritt.
 	\end{itemize}
	\end{bemerkung}
\end{frame}

\begin{frame}{Grenzwertberechnung rekursiv definierter Folgen}
	\footnotesize{
	\begin{example}
		Für die Folge $a_{n}$ mit
		\[
			a_{1} := 0 \qquad a_{n+1} = \sqrt{2 a_{n}}
		\]
		gilt:
		\begin{itemize}
			\item $a_{n}$ konvergiert, weil $a_{n}$ einerseits beschränkt ist, denn $0  \leq a_{1} \leq 2$ und daher für alle $n$
			\[
				0 \leq a_{n+1} = \sqrt{2 a_{n}} = \sqrt{2} \sqrt{a_{n}} \leq \sqrt{2} \sqrt{2} = 2 
			\]
			und andererseits monton wachsend ist:
			\[
				\frac{a_{n+1}}{a_{n}} = \frac{\sqrt{2a_{n}}}{a_{n}} = \sqrt{\frac{2}{a_{n}}} \geq 1
			\]
			\item Die Fixpunktgleichung lautet dann
			\[
				a = \sqrt{2a}
			\] 
			und hat die Lösungen $a=0$ und $a=2$. Da $a_{n}$ monoton wächst und $a_{1} = 1$ ist, muss der Grenzwert $a=2$ sein.
		\end{itemize}
	\end{example}
}
\end{frame}

\subsection{Reihen}

\begin{frame}{Reihen - Definition}
	\begin{definition}
	Für eine beliebige Zahlenfolge $(a_{n})_{n \in \mathbb{N}}$ heißt die Folge der \textbf{Partialsummen} $(s_{n})_{n \in \mathbb{N}}$ mit dem Bildungsgesetz
	\[
		s_{n} := \sum_{i=1}^{n} a_{i}
	\]
	eine \textbf{unendliche Reihe}.
	
	Konvergiert die Folge der Partialsummen $(s_n)$, so heißt die Reihe \textbf{konvergent}. Ansonsten heißt die Reihe \textbf{divergent}.
	
	Konvergiert die Reihe $(s_{n})$, so nennt man ihren Grenzwert 
	\[
		\lim_{n \to \infty} s_{n} = \lim_{n \to \infty} \sum_{i=1}^{n} a_{i} =: \sum_{i=1}^{\infty} a_{i} 
 	\]
 	den \textbf{Wert der Reihe}.
	\end{definition}	
\end{frame}

\begin{frame}{Beispiele für Reihen}
	\begin{example}
		Die \textbf{geometrische Reihe} ist definiert als Folge der Partialsummen der geometrischen Folge 
		\[
			\sum_{i=1}^{n} a_{1} q^{i-1} = a_{1} \sum_{i=1}^{n} q^{i-1} = a_{1} \left( 1 + q + q^{2} + q^{3} + q^{4} + \ldots + q^{n-1}\right)
		\]
	\end{example}
	\begin{example}
		Die Reihe 
		\[
			s_{n}:= \sum_{k=1}^{n} \frac{1}{k(k+1)} = \sum_{k=1}^{n} \left( \frac{1}{k} - \frac{1}{k+1} \right) = 1 - \frac{1}{n+1}
		\] 
		ist offenbar auch explizit als Folge mit $s_{n} = 1-\frac{1}{n+1}$ definiert. An dieser Reihe kann man das wichtige Prinzip \textbf{der Teleskopsummen} erkennen.
	\end{example}
\end{frame}

\begin{frame}{Beispiele für Konvergenz von Reihen}
	\begin{example}
		Die \textbf{harmonische Reihe} divergiert:
		\[
			\sum_{i=1}^{n} \frac{1}{i} = 1 + \frac{1}{2} + \frac{1}{3} + \ldots + \frac{1}{n}
		\] 
		Dagegen konvergiert die \textbf{alternierende harmonische Reihe} 
		\[
			\sum_{i=1}^{n} (-1)^{i+1} \frac{1}{i} = 1 - \frac{1}{2} + \frac{1}{3} - \ldots + (-1)^{n} \frac{1}{n} = \ln 2
		\]
	\end{example}
\end{frame}

\begin{frame}{Konvergenz von Reihen - Notwendige Bedingung}
	\begin{bemerkung}
		\textbf{Notwendig} für die Konvergenz der Reihe $(s_{n})$ ist, dass die Folge der Summanden $(a_{n})$ eine \textbf{Nullfolge} ist, also 
		\[
			\lim_{n \to \infty} s_{n} = \lim_{n \to \infty} \sum_{i=1}^{n} a_{i} \mbox{ konvergiert} \qquad \Rightarrow \lim_{n \to \infty} a_{n} = 0 
		\]
		Diese Bedingung ist aber \textbf{nicht hinreichend}, denn die harmonische Folge 
		\[
			a_{n} = \frac{1}{n} 
		\]
		konvergiert zwar gegen 0, aber die Reihe
		\[
			\sum_{i=1}^{n} a_{i} = \sum_{i=1}^{n} \frac{1}{i} \mbox{ divergiert.}
		\]
	\end{bemerkung}
\end{frame}

\begin{frame}{Arithmetische  Reihe}
	\begin{definition}
		Die Folge der Partialsummen/Teilsummen einer arithmetischen Folge $a_{n} = a_{1} + d (n-1)$, also
		\[
			s_{n} = \sum_{i=1}^{n} a_{i} = \sum_{i=1}^{n} a_{1} + d (i-1)
		\]
		heißt \textbf{arithmetische Reihe} und hat die Summenformel
		\[
			s_{n} = \frac{n}{2}	\left( a_{1} + a_{n} \right) = \frac{n}{2} \left( 2 a_{1} + d (n-1)  \right)
		\]
	\end{definition}
	\begin{bemerkung}
		Die Summenformel hat Carl Friedrich Gauß angeblich als 10-jähriger Schüler gefunden, als er zur Strafe die Zahlen von 1 bis 100 addieren sollte, was eine arithmetische Folge von Zahlen mit $a_{1}=1$ und $d=1$ ist.
	\end{bemerkung}
\end{frame}

\begin{frame}{Arithmetische  Reihe - Gauß}
	\begin{example}
		Die Summenformel hat Gauß so hergeleitet:
		\begin{align*}	
			s_{n} & = a_1  +  & a_{1}  + d   + & \ldots  + & a_{1} + (n-1) d \\
			s_{n} & = a_{1} + (n-1) d  + & a_{1} + (n-2)d    + & \ldots  + & a_{1} \\ 
			\hline
			\Rightarrow 2 s_{n} & = 2 a_{1} + (n-1) d + & 2 a_{1} + (n-1) d + & \ldots + & 2 a_{1} + (n-1) d
		\end{align*}
		Da das $n$-mal derselbe Term ist, ist also
		\[	
			2s_{n} = n \left( 2a_{1} + d(n-1) \right) = n \left( a_{1} + a_{1} + d(n-1) \right) = n \left( a_{1} + a_{n} \right)
		\]
	\end{example}
\end{frame}	

\begin{frame}{Beweis der Summenformel für die arithmetische Reihe}
	\small{	Zu zeigen ist: 
		\[
			s_{n} = \sum_{i=1}^{n} a_{1} + d (i-1) = \frac{n}{2} \left(2a_{1} + d (n-1)\right)
		\]
		Diese Gleichung gilt für $n=1$, denn
		\[
			s_{1} = \sum_{i=1}^{1} a_{1} = \frac{1}{2}\left(2 a_{1} \right)
		\]
		Angenommen nun, die Gleichung gilt für $n$. Wir zeigen nun, dass sie dann auch für $n+1$ gilt:
		\begin{align*}
		& s_{n+1} = \sum_{i=1}^{n+1} a_{i}  = \left( \sum_{i=1}^{n} a_{i} \right) + a_{n+1}  = \frac{n}{2} \left( 2 a_{1} + d(n-1) \right) + a_{1} + d n \\
		& = n a_{1} +  d n - d + a_{1} + dn   = (n+1) a_{1} + 2dn  = \frac{n+1}{2} \left( 2a_{1} + d (n+1 -1) \right)
 		\end{align*}}
\end{frame}
\begin{frame}{Summenformel}
	\begin{example}
		Die Summe der ersten $n$ natürlichen Zahlen ist daher
		\[
			\sum_{i=1}^{n} i = 1 + 2 + 3 + \ldots + n = \frac{n(n+1)}{2}	
	 \]
	 \small{
	 Für die Summe der ersten $n$ ungeraden Zahlen gilt:
	 \[
		 \sum_{i=1}^{n} 1 + 2 \cdot (i-1) = \sum_{i=1}^{n} 2\cdot i -1 =  1 + 3 + 5 + \ldots + 2\cdot n - 1 = \frac{n}{2}\left( 1 + 2n-1 \right) = n^{2}
	 \]
	 Für die Summe der ersten $n$ geraden Zahlen gilt:
	 \[
	 \sum_{i=1}^{n} 2 + 2 \cdot (i-1) = \sum_{i=1}^{n} 2\cdot i =  2 + 4 + 6 + \ldots + 2\cdot n = \frac{n}{2}\left( 2 + 2n \right) = n(n+1)
	 \]
	}
	\end{example}
\end{frame}

\begin{frame}{Geometrische Reihe}
		\begin{definition}
			Die Folge der Partialsummen/Teilsummen einer geometrischen Folge $a_{n} = a_{1} q^{n-1}$, also
			\[
			s_{n} = \sum_{i=1}^{n} a_{i} = \sum_{i=1}^{n} a_{1} q^{i-1}
			\]
			heißt \textbf{geometrische Reihe} und hat für $q\neq 1 $ die Summenformel
			\[
				s_{n} = a_{1} \frac{1-q^{n}}{1-q} = a_{1} \frac{q^{n}-1}{q-1} 
			\]
		\end{definition}	
\end{frame}

\begin{frame}{Herleitung der Summenformel für geometrische Reihen}	
	\begin{example}
		Für die Summenformel gilt wie bei Gauß' Verfahren:
		\begin{align*}	
		s_{n} & = a_1  & + a_{1} q^{1} & + \ldots  & + a_{1}  q^{n-1}  \\
		q s_{n} & = &  a_{1} q^{1} & + \ldots  & + a_{1} q^{n-1} & + a_{1} q^{n}  \\ 
		\hline
		\Rightarrow s_{n} -q s_{n} & = a_{1}  &  &  & & - a_{1} q^{n} 
		\end{align*}
		Es bleiben also zwei Terme übrig:
		\[
			s_{n} \left( 1-q \right) = a_{1} - a_{1} q^{n} = a_{1} \left( 1 - q^{n} \right)
		\]
		und daher (nur) für $q \neq 1$
		\[
			s_{n} = a_{1} \frac{1-q^{n}}{1-q}
		\]	
	\end{example}	
\end{frame}

\begin{frame}{Konvergenz von arithmetischer und geometrischer Reihe}
	\begin{bemerkung}
		Die \textbf{arithmetische Reihe} 
		\[
			\sum_{i=1}^{n} a_{1} + d (i-1) = \frac{n}{1} \left( 2a_{1} + d (n-1) \right) = \frac{(2a_{1}-d)n +d n^{2}}{2}
		\]
		konvergiert offensichtlich nur, wenn sowohl $d=0$ also auch $a_{1}=0$ ist. Deshalb ist die arithmetische Reihe eigentlich immer divergent.
	\end{bemerkung}
	\begin{bemerkung}
		Die \textbf{geometrische Folge} ist nur für $|q|<1$ eine Nullfolge. 
		
		Daher gilt für die \textbf{geometrische Reihe} für $|q|<1$
		\[
		\lim_{n \to \infty} \sum_{i=1}^{n} a_{1} q^{n-1} = \lim_{n \to \infty} a_{1} \frac{1-q^{n}}{1-q} = \frac{a_{1}}{1-q}
		\]
	\end{bemerkung}
\end{frame}

\begin{frame}{Beispiele}
	\begin{example}
		Für $a_{1}=1$ und $q=\frac{3}{4}$ ist
		\[
			\sum_{i=0}^{\infty} \left( \frac{3}{4} \right)^{i} = \lim_{n \to \infty} \sum_{i=0}^{n} \left( \frac{3}{4} \right)^{i} =\lim_{n \to \infty} \frac{1-(\frac{3}{4})^{n+1}}{1-\frac{3}{4}}= \frac{1}{1-\frac{3}{4}} = 4 
		\]
				Für $a_{1}=1$ und $q=\frac{1}{2}$ ist
				\[
				\sum_{i=0}^{\infty} \left( \frac{1}{2} \right)^{i} = \lim_{n \to \infty} \sum_{i=0}^{n} \left( \frac{1}{2} \right)^{i} =\lim_{n \to \infty} \frac{1-(\frac{1}{2})^{n+1}}{1-\frac{1}{2}}= \frac{1}{1-\frac{1}{2}} = 2 
				\]	
	\end{example}
\end{frame}

%\begin{frame}{Rechenregeln für konvergente Reihen}
%	\begin{theorem}
%		Für zwei \textbf{konvergente Reihen}	$\sum_{i=1}^{\infty} a_{i} =: s_{a}$ und  $\sum_{i=1}^{\infty} b_{i} =: s_{b}$ 	gilt: 
%		\begin{itemize}
%			\item Die Reihe der Summen ist konvergent und
%			\[
%				\sum_{i=1}^{\infty} \left( a_{i} + b_{i} \right) = \sum_{i=1}^{\infty}a_{i} + \sum_{i=1}^{\infty} b_{i}   = s_{a} + s_{b} 
%			\] 
%			\item Für jedes $c \in \mathbb{R}$ gilt:
%			\[
%			\sum_{i=1}^{\infty} c a_{i} = c \sum_{i=1}^{\infty} a_{i} = c s_{a}  
%			\]
%			\item Ist $a_{i} \leq b_{i}$ für alle $i \in \mathbb{N}$, so gilt
%			\[
%				\sum_{i=1}^{\infty} a_{i} \leq \sum_{i=1}^{\infty} b_{i}  
%			\]
%		\end{itemize}
%	\end{theorem}
%\end{frame}
%
%\begin{frame}{Absolute Konvergenz}
%	\begin{definition}
%		Eine Reihe $\sum_{i=1}^{\infty} a_{i}$ \textbf{konvergiert absolut}, wenn die Folge der Partialsummen der Beträge der zugehörigen Folge konvergiert:
%		\[
%			\sum_{i=1}^{\infty} | a_{i} | = \lim_{n \to \infty} \sum_{i=1}^{n} | a_{i} | \text{ existiert.}
%		\]
%	\end{definition}
%	\begin{bemerkung}
%		Jede absolut konvergente Reihe konvergiert auch. Die absolute Konvergenz ist also  die stärkere Bedingung, denn jede absolut konvergente Reihe konvergiert auch, aber nicht jede konvergente Reihe konvergiert auch absolut.
%	\end{bemerkung}
%\end{frame}	
%\begin{frame}{Beispiel absolute Konvergenz}
%	\begin{example}
%		Die Reihe 
%		\[
%			\sum\limits_{i=1}^{\infty} \left(-1\right)^{i}\frac{1}{i} = \ln 2
%		\]
%		konvergiert, aber sie konvergiert nicht absolut, denn
%		\[
%			\sum\limits_{i=1}^{\infty} \left| \left(-1 \right)^{i}  \frac{1}{i} \right| = \sum\limits_{i=1}^{\infty} \frac{1}{i} 
%		\]
%		ist die harmonische Reihe und divergiert bekanntlich.
%	\end{example}
%\end{frame}
%
%\subsubsection{Konvergenzkriterien}
%\begin{frame}{Konvergenzkriterien für Reihen}
%	Gegeben sei eine Reihe $\sum_{i=1}^{n} a_{i}$.
%	Dann gilt:
%	\begin{theorem}
%		\begin{itemize}
%			\item \textbf{Nullfolgenkriterium:} Falls die $a_{n}$ \textbf{keine} Nullfolge bilden, \textbf{divergiert} die Reihe $\sum_{i=1}^{n} a_{i}$.
%			\item \textbf{Leibnizkriterium:} Die alternierende Reihe $\sum_{i=1}^{n} (-1)^{i} a_{i}$ \textbf{konvergiert}, falls $a_{i}$ eine \textbf{monoton fallende Nullfolge} ist. Das impliziert $a_{i}\geq 0$ für alle $i$. Für den Wert der Reihe gilt die Abschätzung:
%			\[
%				\left| \sum_{i=1}^{\infty} a_{i} - \sum_{i=1}^{n} a_{i} \right|  = \left| \sum_{i=n+1}^{\infty} a_{i}  \right| \leq a_{n+1} 
%			\] 
%			\item \textbf{Majorantenkriterium:} Die Reihe $\sum_{i=1}^{n} a_{i}$ \textbf{konvergiert absolut}, wenn es eine \textbf{konvergente} Reihe (Majorante) $\sum_{i=1}^{n} b_{i}$ gibt mit
%			\[
%				| a_{i} | \leq b_{i} \qquad \forall i \in \mathbb{N}
%			\]
%		\end{itemize}
%	\end{theorem}
%\end{frame}
%
%\begin{frame}{Konvergenzkriterien für Reihen}
%	\begin{theorem}
%		\begin{itemize}
%			\item \textbf{Minorantenkriterium:} Die Reihe $\sum_{i=1}^{n} a_{i}$ \textbf{divergiert}, falls es eine \textbf{divergente Minorante} $\sum_{i=1}^{n} b_{i}$ gibt mit
%			\[
%				0 \leq b_{i} \leq a_{i} \qquad \forall i \in \mathbb{N}
%			\]
%			\item \textbf{Quotientenkriterium:} Existiert der Grenzwert $r = \lim_{i \to \infty} \left| \frac{a_{i+1}}{a_{i}} \right|$, so gilt:
%			
%				Im Fall $\begin{cases}
%					r < 1 \qquad \text{konvergiert die Reihe (absolut)}\\
%					r >1  \qquad \text{divergiert die Reihe}\\
%					r = 1 \qquad \text{ist alles möglich.}
%				\end{cases}$
%			\item \textbf{Wurzelkriterium:} Existiert der Grenzwert $r = \lim_{i \to \infty} \sqrt[i]{\left| a_{i} \right|}$, so gilt:
%			
%			Im Fall $\begin{cases}
%			r < 1 \qquad \text{konvergiert die Reihe (absolut)}\\
%			r >1  \qquad \text{divergiert die Reihe}\\
%			r = 1 \qquad \text{ist alles möglich.}
%			\end{cases}$
%		\end{itemize}
%	\end{theorem}
%\end{frame}
%
%\begin{frame}{Allgemeine harmonische Reihe}
%	\begin{theorem}
%		Die \textbf{allgemeine harmonische Reihe} mit $\alpha \in \mathbb{R}$
%		\[
%			\sum\limits_{k=1}^{\infty} \frac{1}{k^{\alpha}} 
%		\begin{cases}
%		  	\text{divergiert für } \alpha \leq 1 \\
%			\text{konvergiert für }  \alpha > 1
%		\end{cases}
%		\]
%	\end{theorem}
%	\begin{proof}[Beweis]
%		\scriptsize{
%		Die Reihe $\sum\limits_{k=1}^{\infty} \frac{1}{k}$ divergiert, denn
%		\begin{align*}
%		s_n & = \overset{s_{2^{k+1}}}{\overbrace{1 + \frac{1}{2} + \underset{\geq 2 \cdot \frac{1}{4}}{\underbrace{\frac{1}{3} + \frac{1}{4}}} +\underset{\geq 4 \cdot \frac{1}{8}}{\underbrace{\frac{1}{5} + \frac{1}{6} + \frac{1}{7} + \frac{1}{8}}} + \ldots + \underset{\geq 2^{k} \cdot \frac{1}{2^{k+1}}}{\underbrace{\left( \frac{1}{2^{k} + 1} + \ldots + \frac{1}{2^{k + 1}} \right)}}  } + \ldots + \frac{1}{n} } \\
%		& \geq 1+ \frac{1}{2}  + 2\cdot \frac{1}{4} + 4 \cdot \frac{1}{8} + \ldots + 2^{k} \cdot \frac{1}{2^{k+1}} \geq \frac{3}{2} + k \frac{1}{2} = \frac{k+3}{2}
%		\end{align*}
%		Also gilt für alle $k  \in \mathbb{N}$, dass $s_{2^{k+1}} \geq \frac{k+3}{2} \stackrel{k \to \infty}{\to} + \infty$
%	}
%	\end{proof}
%\end{frame}
%
%\begin{frame}{Fortsetzung des Beweises}
%	\begin{proof}[Beweis]
%	\tiny{Die Divergenz der Reihe für $\alpha \leq 1$ ist damit klar, da die harmonische Reihe, deren Divergenz gerade gezeigt wurde, eine divergente Minorante ist.
%	
%	Für $\alpha > 1$ gilt:
%	\begin{align*}
%	& s_{2^{k}-1}  = 1 + \left( \frac{1}{2^{\alpha}}  +
%	\frac{1}{3^{\alpha}} \right) + \left( \frac{1}{4^{\alpha}} + \frac{1}{5^{\alpha}} + \frac{1}{6^{\alpha}} + \frac{1}{7^{\alpha}} +  \right) + \ldots + \left( \frac{1}{(2^{k-1})^{\alpha}} + \ldots + \frac{1}{(2^{k}-1)^{\alpha}} \right) \\
%	& \leq 1 + \frac{2}{2^{\alpha}} + \frac{4}{4^{\alpha}} + \ldots + \frac{2^{k-1}}{(2^{k-1})^{\alpha}} = \sum\limits_{i=0}^{k-1}\left( \frac{1}{2^{\alpha - 1}} \right)^{i} = \sum\limits_{i=0}^{k-1} q^{i} \text{ mit } q=\frac{1}{2^{\alpha - 1}} 
%	\end{align*}}
%	Für $\alpha > 1$ ist $q<1$ und damit konvergiert die letzte Reihe und ist eine konvergente Majorante für die allgemeine harmonische Reihe mit $\alpha > 1$.
%	\end{proof}
%\end{frame}
%
%\begin{frame}{Beispiele}
%	\begin{itemize}
%		\item Die Reihe $\sum\limits_{i=0}^{\infty} \frac{k^2 +7}{5 k^2 +1 }$ divergiert, da $\frac{k^2 +7}{5 k^2 +1 } \to \frac{1}{5}$ keine Nullfolge ist.
%		\item Die Reihe $\sum\limits_{k=1}^{\infty} \frac{(-1)^{k+1}}{k}$ konvergiert nach dem Leibnizkriterium, weil $\frac{1}{k}$ eine monoton fallende Nullfolge ist.
%		\item Die Reihe $\sum\limits_{k=1}^{\infty} \frac{1}{k^2 -k +1}$ konvergiert, da wegen
%		\[
%			k^2 -k + 1 = (k-1)^2 + k \geq (k-1)^2 \stackrel{k\geq 2}{\Rightarrow} \frac{1}{k^2 -k +1} \leq \frac{1}{(k-1)^2}
%		\] 
%		und wegen der Konvergenz der Reihe $\sum\limits_{k=1}^{\infty} \frac{1}{(k-1)^2}$ eine konvergente Majorante existiert.
%		\item Die Reihe $\sum\limits_{k=0}^{\infty} \frac{2^k}{k!}$ konvergiert wegen des Quotientenkriteriums:
%		\[
%			r = \lim_{k \to \infty} \left| \frac{a_{k+1}}{a_{k}} \right| = \lim_{k \to \infty}\left| \frac{2^{k+1}}{(k+1)!} \frac{k!}{2^k} \right| = \lim_{k \to \infty} \frac{2}{k+1} = 0 < 1
%		\]
%	\end{itemize}
%\end{frame}
%
%\begin{frame}{Beispiele}
%	\begin{itemize}
%		\item Die Reihe $\sum\limits_{k=1}^{\infty} \left( \frac{1}{3} - \frac{1}{\sqrt{k}}\right)^{k}$ konvergiert nach dem Wurzelkriterium:
%		\[
%			r = \lim_{k \to \infty} \sqrt[k]{\left| \left( \frac{1}{3} - \frac{1}{\sqrt{k}} \right)^{k} \right|} = \lim_{k \to \infty} \left| \frac{1}{3} - \frac{1}{\sqrt{k}} \right| = \frac{1}{3} < 1
%		\]
%		\item Die harmonische Reihe $ \sum\limits_{k=1}^{\infty} \frac{1}{k}$ divergiert bekanntlich, für das Quotientenkriterium gilt:
%		\[
%			r = \lim_{k \to \infty} \frac{1}{k+1} k = 1
%		\]
%		\item Die Reihe $ \sum\limits_{k=1}^{\infty} \frac{1}{k^2}$ konvergiert bekanntlich, für das Quotientenkriterium gilt:
%		\[
%		r = \lim_{k \to \infty} \frac{1}{(k+1)^2} k^2 = 1
%		\] 
%	\end{itemize}
%\end{frame}
%
